% GENERATED FILE. Edit canonical lessons, then run npm run generate:books.
% canonical-source-hash: fnv1a64:e54d543c3405d737
% canonical-lessons: UR-C403-sarae, UR-C403-qiyam, UR-C403-thikana, UR-C403-bichhana, UR-C403-lift

\chapter{Inns and Lodgings}
\label{ch:ur-inns-and-lodgings}
\hlchaptermodality{403}

\begin{chapteropening}
\textbf{By the end of this chapter:} \emph{I can find a place to stay for the night and settle in.}

Complete the last lesson of chapter 403: I can find a place to stay for the night and settle in.
\end{chapteropening}

\section[\texorpdfstring{\ur{سرائے} (sarāe)}{سرائے (sarāe)}]{\ur{سرائے} (sarāe) — an inn}

\label{lesson:UR-C403-sarae}



\begin{quote}
Before the new one: say the Urdu for to put (someone) to sleep, then the Urdu for to wake (someone) up.
\end{quote}

\subsection*{You'll want to know: \ur{سرائے}}
\textbf{\ur{سرائے}} — \emph{sarāe} — \textquotedblleft{}an inn\textquotedblright{}.

\begin{quote}
\ur{پرانے} \ur{شہر} \ur{میں} \ur{ایک} \ur{سرائے} \ur{ہے۔}
\end{quote}

\begin{quote}
\emph{purāne shahr meṅ ek sarāe hai.}
\end{quote}

\begin{quote}
There is an inn in the old city.
\end{quote}

Feminine.

\begin{cousinweb}
\textbf{\ur{سرائے}} (\emph{sarāe}) is usually traced to Persian \emph{sarāy}, \textquotedblleft{}a house, a palace\textquotedblright{}; in Urdu it is an inn where travellers stop for the night.
\end{cousinweb}

\begin{grammarlens}[title={What you've built}]
One more word for getting around.
\end{grammarlens}

\subsection*{Your turn}
\begin{itemize}
\raggedright
  \item \emph{Say it:} \emph{sarāe}
  \item \emph{Say it:} \emph{sarāe}, once more
  \item \emph{Say it:} \emph{purāne shahr meṅ ek sarāe hai.} — There is an inn in the old city.
  \item \emph{Recall:} say the Urdu for to put (someone) to sleep, then the Urdu for to wake (someone) up, then say \emph{sarāe} again
\end{itemize}

\begin{tcolorbox}[breakable,colback=teal!4,colframe=teal!35!black,title={Before you move on}]
What does \ur{سرائے} mean? (\textquotedblleft{}An inn\textquotedblright{}.) Say it once more.
\end{tcolorbox}
\section[\texorpdfstring{\ur{قیام} (qiyām)}{قیام (qiyām)}]{\ur{قیام} (qiyām) — a stay}

\label{lesson:UR-C403-qiyam}



\begin{quote}
Before the new one: say the Urdu for to wake (someone) up, then the Urdu for an inn.
\end{quote}

\subsection*{You'll want to know: \ur{قیام}}
\textbf{\ur{قیام}} — \emph{qiyām} — \textquotedblleft{}a stay\textquotedblright{}.

\begin{quote}
\ur{سرائے} \ur{میں} \ur{ہمارا} \ur{قیام} \ur{تین} \ur{دن} \ur{کا} \ur{تھا۔}
\end{quote}

\begin{quote}
\emph{sarāe meṅ hamārā qiyām tīn din kā thā.}
\end{quote}

\begin{quote}
We stayed at the inn for three days.
\end{quote}

Masculine.

\begin{cousinweb}
\textbf{\ur{قیام}} (\emph{qiyām}) is usually traced, through Persian, to Arabic \emph{qiyām}, \textquotedblleft{}standing, remaining\textquotedblright{}; in Urdu it is a stay somewhere.
\end{cousinweb}

\begin{grammarlens}[title={What you've built}]
One more word for getting around.
\end{grammarlens}

\subsection*{Your turn}
\begin{itemize}
\raggedright
  \item \emph{Say it:} \emph{qiyām}
  \item \emph{Say it:} \emph{qiyām}, once more
  \item \emph{Say it:} \emph{sarāe meṅ hamārā qiyām tīn din kā thā.} — We stayed at the inn for three days.
  \item \emph{Recall:} say the Urdu for to wake (someone) up, then the Urdu for an inn, then say \emph{qiyām} again
\end{itemize}

\begin{tcolorbox}[breakable,colback=teal!4,colframe=teal!35!black,title={Before you move on}]
What does \ur{قیام} mean? (\textquotedblleft{}A stay\textquotedblright{}.) Say it once more.
\end{tcolorbox}
\section[\texorpdfstring{\ur{ٹھکانا} (ṭhikānā)}{ٹھکانا (ṭhikānā)}]{\ur{ٹھکانا} (ṭhikānā) — a place to stay; whereabouts}

\label{lesson:UR-C403-thikana}



\begin{quote}
Before the new one: say the Urdu for an inn, then the Urdu for a stay.
\end{quote}

\subsection*{You'll want to know: \ur{ٹھکانا}}
\textbf{\ur{ٹھکانا}} — \emph{ṭhikānā} — \textquotedblleft{}a place to stay; whereabouts\textquotedblright{}.

\begin{quote}
\ur{رات} \ur{کے} \ur{لیے} \ur{ہمیں} \ur{ٹھکانا} \ur{ملا۔}
\end{quote}

\begin{quote}
\emph{rāt ke liye hameṅ ṭhikānā milā.}
\end{quote}

\begin{quote}
We found somewhere to stay for the night.
\end{quote}

Masculine.

\begin{cousinweb}
\textbf{\ur{ٹھکانا}} (\emph{ṭhikānā}) is inherited: CDIAL lists the Hindi cousin \emph{ṭhikānā}, \textquotedblleft{}a station\textquotedblright{}, under \emph{ṭhīkka}, \textquotedblleft{}firm, right\textquotedblright{}, a form it reconstructs (CDIAL 5503); \textbf{\ur{ٹھیک}} comes from the same source.
\end{cousinweb}

\begin{grammarlens}[title={What you've built}]
One more word for getting around.
\end{grammarlens}

\subsection*{Your turn}
\begin{itemize}
\raggedright
  \item \emph{Say it:} \emph{ṭhikānā}
  \item \emph{Say it:} \emph{ṭhikānā}, once more
  \item \emph{Say it:} \emph{rāt ke liye hameṅ ṭhikānā milā.} — We found somewhere to stay for the night.
  \item \emph{Recall:} say the Urdu for an inn, then the Urdu for a stay, then say \emph{ṭhikānā} again
\end{itemize}

\begin{tcolorbox}[breakable,colback=teal!4,colframe=teal!35!black,title={Before you move on}]
What does \ur{ٹھکانا} mean? (\textquotedblleft{}A place to stay; whereabouts\textquotedblright{}.) Say it once more.
\end{tcolorbox}
\section[\texorpdfstring{\ur{بچھانا} (bichhānā)}{بچھانا (bichhānā)}]{\ur{بچھانا} (bichhānā) — to spread out (a sheet); to make up (a bed)}

\label{lesson:UR-C403-bichhana}



\begin{quote}
Before the new one: say the Urdu for a stay, then the Urdu for a place to stay.
\end{quote}

\subsection*{You'll want to know: \ur{بچھانا}}
\textbf{\ur{بچھانا}} — \emph{bichhānā} — \textquotedblleft{}to spread out (a sheet); to make up (a bed)\textquotedblright{}.

\begin{quote}
\ur{سرائے} \ur{والے} \ur{نے} \ur{بستر} \ur{پر} \ur{نئی} \ur{چادر} \ur{بچھائی۔}
\end{quote}

\begin{quote}
\emph{sarāe vāle ne bistar par naī chādar bichhāī.}
\end{quote}

\begin{quote}
The innkeeper spread a fresh sheet on the bed.
\end{quote}

Past, with \ur{نے} (he or she …): \ur{اس} \ur{نے} \ur{چٹائی} \ur{بچھائی}, \emph{us ne chaṭāī bichhāī}.

\begin{cousinweb}
\textbf{\ur{بچھانا}} (\emph{bichhānā}) is inherited: CDIAL lists the Hindi cousin \emph{bichānā}, \textquotedblleft{}to spread\textquotedblright{}, under \emph{vicchādayati}, \textquotedblleft{}spreads out\textquotedblright{}, a form it reconstructs (CDIAL 11692).
\end{cousinweb}

\begin{grammarlens}[title={What you've built}]
One more word for getting around.
\end{grammarlens}

\subsection*{Your turn}
\begin{itemize}
\raggedright
  \item \emph{Say it:} \emph{bichhānā}
  \item \emph{Say it:} \emph{bichhānā}, once more
  \item \emph{Say it:} \emph{sarāe vāle ne bistar par naī chādar bichhāī.} — The innkeeper spread a fresh sheet on the bed.
  \item \emph{Recall:} say the Urdu for a stay, then the Urdu for a place to stay, then say \emph{bichhānā} again
\end{itemize}

\begin{tcolorbox}[breakable,colback=teal!4,colframe=teal!35!black,title={Before you move on}]
What does \ur{بچھانا} mean? (\textquotedblleft{}To spread out (a sheet); to make up (a bed)\textquotedblright{}.) Say it once more.
\end{tcolorbox}
\section[\texorpdfstring{\ur{لفٹ} (lifṭ)}{لفٹ (lifṭ)}]{\ur{لفٹ} (lifṭ) — a lift}

\label{lesson:UR-C403-lift}



\begin{quote}
Before the new one: say the Urdu for a place to stay, then the Urdu for to spread out (a sheet).
\end{quote}

\subsection*{You'll want to know: \ur{لفٹ}}
\textbf{\ur{لفٹ}} — \emph{lifṭ} — \textquotedblleft{}a lift\textquotedblright{}.

\begin{quote}
\ur{ہوٹل} \ur{کی} \ur{لفٹ} \ur{بند} \ur{تھی۔}
\end{quote}

\begin{quote}
\emph{hoṭal kī lifṭ band thī.}
\end{quote}

\begin{quote}
The hotel lift was out of service.
\end{quote}

Feminine.

\begin{cousinweb}
\textbf{\ur{لفٹ}} (\emph{lifṭ}) is borrowed from English \emph{lift}.
\end{cousinweb}

\begin{grammarlens}[title={What you've built}]
One more word for getting around.
\end{grammarlens}

\subsection*{Your turn}
\begin{itemize}
\raggedright
  \item \emph{Say it:} \emph{lifṭ}
  \item \emph{Say it:} \emph{lifṭ}, once more
  \item \emph{Say it:} \emph{hoṭal kī lifṭ band thī.} — The hotel lift was out of service.
  \item \emph{Recall:} say the Urdu for a place to stay, then the Urdu for to spread out (a sheet), then say \emph{lifṭ} again
\end{itemize}

\begin{tcolorbox}[breakable,colback=teal!4,colframe=teal!35!black,title={Before you move on}]
What does \ur{لفٹ} mean? (\textquotedblleft{}A lift\textquotedblright{}.) Say it once more.
\end{tcolorbox}
